The bifurcation surface named below lies beyond one flat chart, and the statements made there hold on the cross-sections of the horizon locus $r = 2M$ in the Painlev\'e-Gullstrand chart, where Propositions \claimref{singularity.unipotent-ray} and \claimref{singularity.solution-where-divergence-lies} are stated (read through Fact \claimref{synthesis.imported-gtg-dictionary}). The rate statements are the cited propositions (the divergent rate is \claimref{singularity.connection-bivector-computed}(3)). For the locus: the fixed point set of the Killing boost is where the boost Killing field vanishes, which for Schwarzschild is the bifurcation surface $r = 2M$; the Euclidean section of Proposition \claimref{singularity.euclidean-section-period}, whose period Theorem \claimref{singularity.rate-discharge-hypothesis-iii} uses, is the region $r \geq 2M$ and does not contain the singular boundary at all. Proposition \claimref{singularity.unipotent-ray} shows $\hb$ invertible in closed form there and Proposition \claimref{singularity.solution-where-divergence-lies}(3) shows $\det\hb = 1$, so the chain's own Step 1 (locate the failure, Step 1 of \claimref{transport.transport-algorithm}) classifies the horizon as a locus where nothing algebraic happens. A premise imposed there is therefore outside the principle's scope by the principle's own diagnosis.
